\documentclass[12pt]{article}

\usepackage[margin=1in]{geometry}
\usepackage{parskip}
\usepackage{amsmath,amssymb}

\begin{document}

\begin{center}
  {\large\bfseries From Interpretable to Black-Box:\\Dynamic Model Identification
  for Small USV Control Design}\\[0.8em]
  Brian S. Bingham and Carson Vogt\\
  Department of Mechanical and Aerospace Engineering\\
  Naval Postgraduate School, Monterey, CA, USA\\
  \texttt{bbingham@nps.edu}, \texttt{crvogt@nps.edu}
\end{center}

\bigskip

Effective low-level feedback control of small underactuated uncrewed surface
vehicles (USVs), specifically vessels under seven meters with a single thruster
and rudder, requires dynamic models that provide both actionable design insight and
predictive accuracy. Tow-tank testing is impractical at this scale;
parameters must be identified solely from limited field trials. Existing system
identification methods span a wide range of model complexity but offer little
guidance on how much fidelity is enough, because sufficiency is only meaningful
relative to a stated purpose. Absent that specification, the default is a bias
toward more fidelity, since the benefit of an added term is readily quantified
while its costs in identifiability, experimental burden and obscured design
insight are not.

We propose a methodology built around a deliberate spectrum of three model levels,
each identified from the same experimental data: (1) first-order, linear,
single-input single-output (SISO) transfer functions, Nomoto-style models for
surge and yaw that are immediately interpretable and support controller architecture
intuition, but risk oversimplification when considered in isolation;
(2) a physics-based nonlinear state-space formulation incorporating quadratic
hydrodynamic drag; and (3) a Nonlinear AutoRegressive with eXogenous inputs (NARX)
multilayer perceptron. Rather than seeking a single best model, the spectrum is
used collectively: simple models anchor human intuition; structured nonlinear models
expose which terms are dynamically salient; and the black-box model bounds the
benefit of added complexity.

Preliminary results from step-response field experiments demonstrate the approach.
First-order SISO models achieve $R^2=0.93$ (surge, $\tau=1.95$ s) and $R^2=0.83$
(yaw rate, $\tau=0.47$ s), providing direct controller tuning parameters.
Cross-model diagnostics identify quadratic drag as the dominant surge nonlinearity
and reveal that speed-dependent rudder effectiveness, a key coupling for
underactuated control design, is not resolvable from a single short trial,
motivating targeted data collection. A margin-based diagnostic makes the
sufficiency question computable, reporting the forward speed at which the
first-order model stops changing the control design. The NARX model matches one-step prediction
accuracy ($R^2=0.93$) but diverges in closed-loop simulation, demonstrating that
physical structure provides essential regularization under limited data. A
larger dataset spanning multiple operating regimes is being collected to
evaluate the full methodology.

\end{document}
